%% Dans ce fichier, les �ventuels accents sont cod�s
%% en ISO-8859-15 (aka latin9)
%% La ligne suivante donne trois exemples 
%% e accent aigu : �, e accent grave : �, e-dans-l'o : �
%%%%%ENTREES%%%%
\begin{tabarg}
\multicolumn{3}{l}{\titretab {Entr\'ees}}\\
\hline
$X$ & Vecteur, \texttt{float[$N_{level}$]} & Signal lidar\\
$Z$ & Vecteur, \texttt{float[$N_{level}$]} & Altitudes r�elles des portes du lidar ($m$)\\
$T_0$ & Coeff., \texttt{float} & Temp�rature � Z=0 ($K$ ou $Celsius$)\\
$P_0$ & Coeff., \texttt{float} & Pression � Z=0 ($hPa$)\\
$\lambda$ & Coeff., \texttt{float} & Longueur d'onde du signal lidar ($nm$)\\
\end{tabarg}

\medskip
%%%%%SORTIES%%%%
\begin{tabarg}
\multicolumn{3}{l}{\titretab {Sorties}}\\
\hline
$\alpha_{mol}$ & Vecteur, \texttt{float[$N_{level}$]} & Profil du coefficient d'extinction mol�culaire ($m^{-1}$)\\
$\beta_{mol}$ & Vecteur, \texttt{float[$N_{level}$]} &  Profil du coefficient de r�tro-diffusion mol�culaire ($m^{-1}$)\\
$T_{rayleigh}$ & Vecteur, \texttt{float[$N_{level}$]} & Profil du coefficient de Rayleigh \\
$Pr^2_{mol}$ & Vecteur, \texttt{float[$N_{level}$]} & Signal mol�culaire r�tro-diffus� ($Photons.m^2$)\\
\end{tabarg}

